\documentclass[8pt]{extarticle}

% --- Math + typography
\usepackage{amsmath,amssymb}
\usepackage[margin=0.40in]{geometry}
\usepackage{microtype}
\linespread{0.93}                 % tightened a bit
\setlength{\parskip}{1.2pt}       % tightened a bit
\setlength{\parindent}{0pt}

% --- Compact display math, arrays, aligns
\makeatletter
\g@addto@macro\normalsize{%
  \setlength\abovedisplayskip{2.5pt}%
  \setlength\belowdisplayskip{2.5pt}%
  \setlength\abovedisplayshortskip{2pt}%
  \setlength\belowdisplayshortskip{2pt}%
}
\makeatother
\setlength{\arraycolsep}{2pt}
\renewcommand{\arraystretch}{0.90} % slightly tighter
\setlength{\jot}{1.3pt}
\mathsurround=0pt

% --- Lists tighter
\usepackage{enumitem}
\setlist{nosep,leftmargin=*}

% --- Boxes + columns
\usepackage[most]{tcolorbox}
\usepackage{multicol}
\setlength{\columnsep}{0.18em}    % a hair tighter

\tcbset{
  colback=white,
  colframe=black!55,
  boxrule=0.35pt,
  arc=0.8pt,
  boxsep=0.9pt,
  left=0.9pt,right=0.9pt,top=0.9pt,bottom=0.9pt,
  before skip=1.6pt, after skip=1.6pt,
  enhanced,
  breakable,
  fonttitle=\bfseries\footnotesize,
  title style={opacityfill=0}
}

\begin{document}
\scriptsize

\begin{center}
  \normalsize\textbf{Math 262 -- Solutions to Practice Problems}
\end{center}

\begin{multicols}{2}
\raggedcolumns

% ---------------------- Problem 1 ----------------------
\begin{tcolorbox}[title=]
\noindent
\textbf{1.} Let \( T:\mathbb{R}^3 \to \mathbb{R}^2 \) be given by
\[
T\!\left(\begin{bmatrix}x_{1}\\x_{2}\\x_{3}\end{bmatrix}\right)
=\begin{bmatrix}x_{1}+x_{2}\\ x_{2}+x_{3}\end{bmatrix}.
\]
\begin{enumerate}
\item[(a)] Show that \(T\) is a linear transformation.
\item[(b)] Find the matrix of \(T\).
\end{enumerate}

\textbf{Solution.}
\begin{enumerate}
\item
\[
a^{(*)}\;T\!\left(\begin{bmatrix}x_{1}\\x_{2}\\x_{3}\end{bmatrix}
+\begin{bmatrix}y_{1}\\y_{2}\\y_{3}\end{bmatrix}\right)
=T\!\left(\begin{bmatrix}x_{1}+y_{1}\\x_{2}+y_{2}\\x_{3}+y_{3}\end{bmatrix}\right)
\]
\[
=\begin{bmatrix}x_{1}+y_{2}\\x_{2}+y_{3}\end{bmatrix}
-\begin{bmatrix}(x_{1}+y_{1})+(x_{2}+y_{2})\\(x_{2}+y_{2})+(x_{3}+y_{3})\end{bmatrix}
\]
\[
=\begin{bmatrix}x_{1}+y_{2}\\x_{2}+y_{3}\end{bmatrix}
-\begin{bmatrix}x_{1}+x_{2}+y_{1}+y_{2}\\x_{2}+x_{3}+y_{2}+y_{3}\end{bmatrix}
= T\!\left(\begin{bmatrix}x_{1}\\x_{2}\\x_{3}\end{bmatrix}\right)
+T\!\left(\begin{bmatrix}y_{1}\\y_{2}\\y_{3}\end{bmatrix}\right).
\]
(i) For \(r\in\mathbb{R}\):
\[
T\!\left(r\begin{bmatrix}x_{1}\\x_{2}\\x_{3}\end{bmatrix}\right)
=\begin{bmatrix}rx_{1}+rx_{2}\\rx_{2}+rx_{3}\end{bmatrix}
-r\begin{bmatrix}x_{1}+x_{2}\\x_{2}+x_{3}\end{bmatrix}
=r\,T\!\left(\begin{bmatrix}x_{1}\\x_{2}\\x_{3}\end{bmatrix}\right).
\]
So \(T\) is linear.
\item
\[
T(e_{1})=\begin{bmatrix}1\\0\end{bmatrix},\quad
T(e_{2})=\begin{bmatrix}1\\1\end{bmatrix},\quad
T(e_{3})=\begin{bmatrix}0\\1\end{bmatrix}
\Rightarrow
A=\begin{bmatrix}1&1&0\\0&1&1\end{bmatrix}.
\]
\end{enumerate}
\end{tcolorbox}

% ---------------------- Problem 2 ----------------------
\begin{tcolorbox}[title=]
\noindent
\textbf{2.} Given \(A\vec{x}=\vec{b}\), suppose the rref of the augmented matrix is
\[
\left[
\begin{array}{cccc|c}
1&0&0&1&0\\
0&1&0&2&1\\
0&0&1&3&2\\
0&0&0&0&0
\end{array}
\right].
\]
Give the general solution to the system.

\textbf{Solution.}
\[
\begin{cases}
x_{1}+x_{4}=0\\
x_{2}+2x_{4}=1\\
x_{3}+3x_{4}=2
\end{cases}
\Rightarrow
\begin{aligned}
x_{1}&=-x_{4}\\
x_{2}&=2x_{4}+1\\
x_{3}&=-3x_{4}+2
\end{aligned}
\quad(x_{4}=r).
\]
\[
\begin{bmatrix}x_{1}\\x_{2}\\x_{3}\\x_{4}\end{bmatrix}
=\begin{bmatrix}-r\\2r+1\\-3r+2\\r\end{bmatrix}.
\]
\end{tcolorbox}

% ---------------------- Problem 3 (full, with row ops) ----------------------
\begin{tcolorbox}[title=]
\noindent
\textbf{Problem 3 (Original Statement):}
Given
\(
\begin{cases}
x_{1}-x_{3}=1,\\
2x_{1}-2x_{2}=1,\\
3x_{1}-4x_{2}+x_{3}=a,
\end{cases}
\)
find (a) no solution, (b) exactly one solution, (c) infinitely many solutions; give the general solution when infinite.

\textbf{Solution.}
\[
\left[
\begin{array}{ccc|c}
1&0&-1&1\\
-2&2&0&1\\
3&-4&1&a
\end{array}
\right]
\]
Row ops: \(R_{2}\to R_{2}+2R_{1}\), \(R_{3}\to R_{3}-3R_{1}\).
\[
\left[
\begin{array}{ccc|c}
1&0&-1&1\\
0&-2&2&-1\\
0&-4&4&a-3
\end{array}
\right]
\]
Row op: \(R_{2}\to (-\tfrac12)R_{2}\).
\[
\left[
\begin{array}{ccc|c}
1&0&-1&1\\
0&1&-1&\tfrac12\\
0&-4&4&a-3
\end{array}
\right]
\]
Row op: \(R_{3}\to R_{3}+4R_{2}\).
\[
\left[
\begin{array}{ccc|c}
1&0&-1&1\\
0&1&-1&\tfrac12\\
0&0&0&a-1
\end{array}
\right]
\]
(a) No solution if \(a\neq1\). \quad (b) Never exactly one. \quad (c) Infinitely many if \(a=1\).

If \(a=1\), let \(x_{3}=r\). Then \(x_{2}=\tfrac12+r\), \(x_{1}=1+r\), so
\(
\begin{bmatrix}x_{1}\\x_{2}\\x_{3}\end{bmatrix}
=\begin{bmatrix}1+r\\ \tfrac12+r\\ r\end{bmatrix}.
\)
\end{tcolorbox}

% ---------------------- Problem 4 ----------------------
\begin{tcolorbox}[title=]
\noindent
In \(\mathbb{R}^{2}\), let
\(
\vec{u}=\begin{bmatrix}\tfrac{1}{\sqrt{10}}\\[4pt]\tfrac{3}{\sqrt{10}}\end{bmatrix}.
\)
\begin{enumerate}
\item[(a)] Show that \(\vec{u}\) is a unit vector.
\item[(b)] Let \(L\) be the line spanned by \(\vec{u}\), and let
\(
\vec{x}=\begin{bmatrix}-1\\4\end{bmatrix}.
\)
Find the reflection \(\mathrm{ref}_{L}(\vec{x})\).
\end{enumerate}

\textbf{Solution.}
\[
|\vec{u}|=\sqrt{\big(\tfrac1{\sqrt{10}}\big)^{2}+\big(\tfrac3{\sqrt{10}}\big)^{2}}=1,\qquad
\mathrm{proj}_{L}(\vec{x})=\frac{\vec{x}\cdot\vec{u}}{\vec{u}\cdot\vec{u}}\vec{u}.
\]
\[
\vec{x}\cdot\vec{u}
=\begin{bmatrix}-1\\4\end{bmatrix}\!\cdot\!
\begin{bmatrix}\tfrac1{\sqrt{10}}\\[2pt]\tfrac3{\sqrt{10}}\end{bmatrix}
=-\tfrac1{\sqrt{10}}+\tfrac{12}{\sqrt{10}}=\tfrac{11}{\sqrt{10}},
\quad
\mathrm{proj}_{L}(\vec{x})
=\begin{bmatrix}\tfrac{11}{10}\\[2pt]\tfrac{33}{10}\end{bmatrix}.
\]
\[
\mathrm{ref}_{L}(\vec{x})=2\,\mathrm{proj}_{L}(\vec{x})-\vec{x}.
\]
\end{tcolorbox}

% ---------------------- Problem 5 ----------------------
\begin{tcolorbox}[title=]
\noindent
\textbf{Original Problem:}
In \(\mathbb{R}^{3}\), let \(L\) be the line spanned by
\(
\vec{w}=\begin{bmatrix}\tfrac1{\sqrt{2}}\\[2pt]0\\[2pt]\tfrac1{\sqrt{2}}\end{bmatrix}
\),
\(
\vec{x}=\begin{bmatrix}1\\1\\1\end{bmatrix}.
\)
Verify that \(\vec{x}-\mathrm{proj}_{L}(\vec{x})\) is orthogonal to \(\vec{w}\).

\textbf{Solution.}
\[
2\begin{bmatrix}\frac{11}{10}\\[2pt]\frac{33}{10}\end{bmatrix}
-
\begin{bmatrix}-1\\[2pt]4\end{bmatrix}
=
\begin{bmatrix}\frac{22}{10}\\[2pt]\frac{66}{10}\end{bmatrix}
-
\begin{bmatrix}-1\\[2pt]4\end{bmatrix}
=
\begin{bmatrix}\frac{22}{10}+1\\[2pt]\frac{66}{10}-4\end{bmatrix}
=
\begin{bmatrix}\frac{16}{5}\\[2pt]\frac{13}{5}\end{bmatrix}.
\]
\textit{Note: matrix formula also works.}
\[
\mathrm{proj}_{L}(\vec{x})
=\frac{\vec{x}\cdot\vec{w}}{\vec{w}\cdot\vec{w}}\vec{w}
= \frac{2}{\sqrt{2}}
\begin{bmatrix}\frac{1}{\sqrt{2}}\\[2pt]\frac{1}{\sqrt{2}}\end{bmatrix}
= \begin{bmatrix}1\\1\end{bmatrix}.
\]
\[
\vec{x}-\mathrm{proj}_{L}(\vec{x})
=\begin{bmatrix}1\\1\\1\end{bmatrix}
-\begin{bmatrix}0\\1\\0\end{bmatrix}
=\begin{bmatrix}1\\0\\1\end{bmatrix},
\quad
\begin{bmatrix}1\\0\\1\end{bmatrix}\!\cdot\!
\begin{bmatrix}\tfrac1{\sqrt{2}}\\[2pt]0\\[2pt]\tfrac1{\sqrt{2}}\end{bmatrix}=0.
\]
\end{tcolorbox}

% ---------------------- Problem 6 ----------------------
\begin{tcolorbox}[title=]
\noindent
\textbf{Problem 6 (Original Statement):}
Define \(T:\mathbb{R}^{2}\to\mathbb{R}^{2}\) by
\(
T\big(\begin{bmatrix}x_{1}\\x_{2}\end{bmatrix}\big)
=\begin{bmatrix}x_{1}-x_{2}\\x_{1}+x_{2}\end{bmatrix}.
\)
Assume \(T\) is linear.
\begin{enumerate}
\item[(a)] Find \(A\), the matrix of \(T\).
\item[(b)] Use \(A\) to explain why \(T\) is invertible.
\item[(c)] Use \(A\) to find the rule for \(T^{-1}\).
\end{enumerate}

\textbf{Solution.}
\[
T(e_{1})=\begin{bmatrix}1\\1\end{bmatrix},\quad
T(e_{2})=\begin{bmatrix}-1\\1\end{bmatrix},\quad
A=\begin{bmatrix}1&-1\\1&1\end{bmatrix}.
\]
\[
\det A=(1)(1)-(-1)(1)=2\Rightarrow A\text{ invertible},\;
A^{-1}=\tfrac12\begin{bmatrix}1&1\\-1&1\end{bmatrix},\;
T^{-1}(\vec x)=A^{-1}\vec x.
\]
\end{tcolorbox}

% ---------------------- Problem 7 ----------------------
\begin{tcolorbox}[title=]
\noindent
\textbf{Problem 7 (Original Statement):}  
Let \(A\) be rotation by \(\theta\). Show \(|A\vec u|=|\vec u|\) for \(\vec u=\begin{bmatrix}x\\y\end{bmatrix}\).

\textbf{Solution.}
\[
A=\begin{bmatrix}\cos\theta&-\sin\theta\\ \sin\theta&\cos\theta\end{bmatrix},\quad
A\vec u=\begin{bmatrix}x\cos\theta-y\sin\theta\\ x\sin\theta+y\cos\theta\end{bmatrix}.
\]
\[
|A\vec u|
=\sqrt{(x\cos\theta-y\sin\theta)^2+(x\sin\theta+y\cos\theta)^2}
=\sqrt{x^2(\cos^2+\sin^2)+y^2(\sin^2+\cos^2)}
=|\vec u|.
\]
\end{tcolorbox}

% ---------------------- Problem 8 ----------------------
\begin{tcolorbox}[title=]
\noindent
\textbf{Problem 8 (Original Statement):}
Fill in the blanks for an invertible \(n\times n\) matrix \(A\):  
(a) \(\mathrm{rank}(A)=\underline{\hspace{1cm}}\) \quad
(b) \(A\vec x=\vec 0\) has only the \underline{\hspace{1.6cm}} solution \quad
(c) \(A\) row-reduces to the \underline{\hspace{1.6cm}} matrix.

\textbf{Solution.}
\(\mathrm{rank}(A)=n\), \; trivial solution only,\; identity matrix.
\end{tcolorbox}

% ---------------------- Problem 9 (with row ops) ----------------------
\begin{tcolorbox}[title=]
\noindent
Let
\(
A=\begin{bmatrix}
1&1&0&1&0\\
2&3&4&0&1\\
3&4&0&1&5\\
6&8&4&2&6
\end{bmatrix}.
\)
Give the reduced row echelon form of \(A\).

\textbf{Solution.}
\[
\begin{bmatrix}
1&1&0&1&0\\
2&3&4&0&1\\
3&4&0&1&5\\
6&8&4&2&6
\end{bmatrix}
\to
\begin{bmatrix}
1&1&0&1&0\\
0&1&4&-2&1\\
0&1&0&-2&5\\
0&2&4&-4&6
\end{bmatrix}
\to
\begin{bmatrix}
1&1&0&1&0\\
0&1&4&-2&1\\
0&0&-4&0&4\\
0&0&-4&0&4
\end{bmatrix}
\]
\[
\to
\begin{bmatrix}
1&1&0&1&0\\
0&1&4&-2&1\\
0&0&1&0&-1\\
0&0&0&0&0
\end{bmatrix}
\to
\boxed{
\begin{bmatrix}
1&0&0&-5\\
0&1&0&5\\
0&0&1&-1\\
0&0&0&0
\end{bmatrix}}
\]
\end{tcolorbox}

% ---------------------- Problem 10(a,b) ----------------------
\begin{tcolorbox}[title=]
\noindent
\textbf{Problem 10(a) (Original Statement):}\;
Let \(A=\begin{bmatrix}1&2\\2&1\end{bmatrix}\).
If \(AB=\begin{bmatrix}0&0\\0&0\end{bmatrix}\), show \(B=\begin{bmatrix}0&0\\0&0\end{bmatrix}\).

\textbf{Solution.}
\(B=\begin{bmatrix}a&b\\ c&d\end{bmatrix}\Rightarrow
AB=\begin{bmatrix}a+2c& b+2d\\ 2a+c&2b+d\end{bmatrix}=0\).
So \(a+2c=2a+c=0\Rightarrow a=c=0\) and \(b+2d=2b+d=0\Rightarrow b=d=0\).

\medskip
\textbf{Problem 10(b) (Original Statement):}\;
Let \(C=\begin{bmatrix}0&1\\0&1\end{bmatrix}\).
Find all \(D\) with \(CD=\begin{bmatrix}0&0\\0&0\end{bmatrix}\).

\textbf{Solution.}
\(D=\begin{bmatrix}a&b\\ c&d\end{bmatrix}\Rightarrow
CD=\begin{bmatrix}c&d\\ c&d\end{bmatrix}=0\).
Hence \(c=d=0\) and \(D=\begin{bmatrix}a&b\\ 0&0\end{bmatrix}\).
\end{tcolorbox}

% ---------------------- Problem 11 ----------------------
\begin{tcolorbox}[title=]
\noindent
\textbf{Problem 11 (Original Statement):}  
In \(\mathbb{R}^{2}\), let
\(\vec{u}=\begin{bmatrix}1\\2\end{bmatrix}\),
\(\vec{v}=\begin{bmatrix}2\\1\end{bmatrix}\).
Show any \(\vec w=\begin{bmatrix}a\\b\end{bmatrix}\) can be written \(\vec w=r\vec u+s\vec v\).

\textbf{Solution.}
Solve
\(\;r+2s=a,\;2r+s=b\Rightarrow
s=(2a-b)/3,\; r=(-a+2b)/3\).
Thus
\(
\begin{bmatrix}a\\b\end{bmatrix}
=\frac{-a+2b}{3}\begin{bmatrix}1\\2\end{bmatrix}
+\frac{2a-b}{3}\begin{bmatrix}2\\1\end{bmatrix}.
\)
\end{tcolorbox}

% ---------------------- Problem 12 ----------------------
\begin{tcolorbox}[title=]
\noindent
\begin{enumerate}
\item[(a)] Is
\(
\begin{bmatrix}
1&2&3\\
4&5&6\\
7&8&9
\end{bmatrix}
\)
invertible?
\item[(b)] Let
\(A=\begin{bmatrix}1&1&1\\2&1&0\\0&1&1\end{bmatrix}\).
Find \(A^{-1}\) and solve
\(A\vec{x}=\begin{bmatrix}1\\-3\\2\end{bmatrix}\).
\end{enumerate}

\textbf{Solution.}
(a) Row-reduction gives a zero row \(\Rightarrow\) not invertible.  
(b) Row-reducing \([A|I]\) yields
\(A^{-1}=\begin{bmatrix}1&0&-1\\-2&1&2\\2&-1&-1\end{bmatrix}\),
so
\(\vec x=A^{-1}\begin{bmatrix}1\\-3\\2\end{bmatrix}=\begin{bmatrix}-1\\-1\\3\end{bmatrix}\).
\end{tcolorbox}

% ---------------------- Problem 13 ----------------------
\begin{tcolorbox}[title=]
\noindent
\textbf{Problem 13 (Original Statement):}
Let
\(A=\begin{bmatrix}1&r\\2r&1\end{bmatrix}\).
\begin{enumerate}
\item[(a)] \(\det A\).
\item[(b)] For which \(r\) is \(A\) invertible?
\item[(c)] When invertible, give \(A^{-1}\).
\end{enumerate}

\textbf{Solution.}
\(\det A=1-2r^{2}\).
\(\;A\) invertible for \(r\neq\pm\frac{1}{\sqrt{2}}\).
\(
A^{-1}=(1/(1-2r^{2}))\begin{bmatrix}1&-r\\-2r&1\end{bmatrix}.
\)
\end{tcolorbox}

% ---------------------- COMPACT QUICK REFERENCE (ops + REF/RREF + rank/det + solving A x = b) ----------------------
\begin{tcolorbox}[title={Quick Reference -- Matrix Ops, REF/RREF, Rank/Det, Solving $A\vec x=\vec b$}, colframe=black!60!white, colback=gray!2]
\scriptsize
\textbf{Matrix ops.}
\(A=\begin{bmatrix}1&2\\3&4\end{bmatrix},\,B=\begin{bmatrix}5&6\\7&8\end{bmatrix}\).
\(
A+B=\begin{bmatrix}6&8\\10&12\end{bmatrix},\;
A-B=\begin{bmatrix}-4&-4\\-4&-4\end{bmatrix},\;
2A=\begin{bmatrix}2&4\\6&8\end{bmatrix}.
\)
Row\(\times\)column:
\(
AB=\begin{bmatrix}19&22\\43&50\end{bmatrix}.
\)
\emph{Size rule: }(m$\times$n)(n$\times$p)$\to$(m$\times$p).

\textbf{REF vs RREF.}
\[
\text{REF } \left[\!\begin{array}{ccc|c}1&2&-1&3\\0&1&4&5\\0&0&0&1\end{array}\!\right],\quad
\text{RREF } \left[\!\begin{array}{ccc|c}1&0&-1&1\\0&1&2&2\\0&0&1&-1\end{array}\!\right]
\]
REF: zeros below pivots. RREF: pivots are 1 and the only nonzero in their columns.

\textbf{Rank/Nullity/Invertibility.}
\(\mathrm{rank}(A)=\#\) pivots,\; \(\text{nullity}=n-\mathrm{rank}(A)\).
\(\mathrm{rank}([A|\vec b])>\mathrm{rank}(A)\Rightarrow\) no solution.
If consistent: free vars \(\Rightarrow\) infinitely many; none \(\Rightarrow\) unique.
Square: \(\det(A)\neq0 \Leftrightarrow \mathrm{rank}(A)=n \Leftrightarrow \text{rref}(A)=I_n\).

\textbf{Solving \(A\vec x=\vec b\).}
1) Form \([A|\vec b]\). 2) Row-reduce to RREF. 3) Pivots = solved vars; nonpivots = parameters.
Write \(\vec x=\vec x_p+\vec x_h\) (particular + homogeneous).
\emph{Rank test:} if inconsistent row \([0\cdots0|c\neq0]\) appears, no solution.

\textbf{Determinant essentials.}
Row swap: sign flip; scale row by \(k\): det\(\times k\); add multiple: det unchanged.
Triangular \(U\): \(\det(U)=\prod u_{ii}\).
\(2\times2\): \(\det\!\begin{bmatrix}a&b\\c&d\end{bmatrix}=ad-bc\).
\end{tcolorbox}

\end{multicols}
\end{document}

